\section{Results}
The surface morphologies and porous characteristics of the ETLs were characterized using nitrogen isotherms. The synapse-engineered mesoporous $\text{TiO}_2$ exhibited a BET surface area of $142\,\text{m}^2/\text{g}$, representing a massive increase over the $85\,\text{m}^2/\text{g}$ observed for the baseline sample. This expanded surface area provides abundant active sites for high-density sensitization. As shown in the cell schematic (Figure~\ref{fig:schematic}), the synapse-engineered ETL supports a significantly higher loading density of CdSe quantum dots, which are distributed uniformly throughout the mesopores, in contrast to the disordered baseline matrix.

\begin{figure*}[t]
  \centering
  \begin{tikzpicture}[
    scale=0.95,
    every node/.style={font=\sffamily\small},
    layer/.style={draw=primary!80, thick, rounded corners=2pt, align=center},
    qd/.style={circle, draw=accent!80, fill=accent!40, inner sep=0pt, minimum size=6pt}
  ]
    % Draw layers
    % FTO Glass
    \draw[layer, fill=primary!10] (0,0) rectangle (14,0.8) node[midway] {\textbf{FTO Glass Substrate} (Fluorine-doped Tin Oxide)};
    
    % TiO2 ETL
    \draw[layer, fill=secondary!15, draw=secondary] (0,1.0) rectangle (8,2.8) node[midway, align=center] {\textbf{Synapse-Engineered TiO$_2$ ETL}\\Mesoporous Framework ($142\,\text{m}^2/\text{g}$)};
    
    % Baseline TiO2
    \draw[layer, fill=textmute!10, draw=textmute] (9,1.0) rectangle (14,2.8) node[midway, align=center] {\textbf{Baseline TiO$_2$ ETL}\\Disordered Network};
    
    % CdSe Quantum Dots on Synapse-Engineered
    \foreach \x in {0.5, 1.2, 1.8, 2.5, 3.1, 3.8, 4.5, 5.2, 5.9, 6.6, 7.3, 7.8} {
      \node[qd] at (\x, 1.2) {};
      \node[qd] at (\x+0.2, 1.8) {};
      \node[qd] at (\x-0.1, 2.3) {};
      \node[qd] at (\x+0.1, 2.6) {};
    }
    
    % CdSe Quantum Dots on Baseline
    \foreach \x in {9.5, 10.4, 11.2, 12.1, 12.9, 13.6} {
      \node[qd] at (\x, 1.3) {};
      \node[qd] at (\x+0.1, 2.0) {};
      \node[qd] at (\x-0.1, 2.5) {};
    }
    
    % Electrolyte interface
    \draw[dashed, thick, color=accent] (-0.5,3.2) -- (14.5,3.2) node[above, midway] {Polysulfide Electrolyte ($S^{2-}/S_x^{2-}$ Interface)};
    
    % Counter Electrode
    \draw[layer, fill=primary!20] (0,3.8) rectangle (14,4.6) node[midway] {\textbf{Counter Electrode} (Cu$_2$S / Brass)};
    
    % Light source indicator
    \draw[ultra thick, orange, ->] (3, 5.6) -- (3, 4.8) node[above, at start, color=orange!80!black] {Incident Light ($h\nu$)};
    \draw[ultra thick, orange, ->] (11, 5.6) -- (11, 4.8);
    
    % Annotation boxes
    \node[anchor=west, text=primary] at (0.2, 3.4) {\faArrowRight\ High Active Site Density};
    \node[anchor=west, text=textmute] at (9.2, 3.4) {\faTimes\ Lower QD Loading};
    
  \end{tikzpicture}
  \caption{Schematic comparison of the synapse-engineered mesoporous TiO$_2$ electron transport layer (left) displaying high quantum dot loading and uniform distribution, compared to the disordered baseline TiO$_2$ framework (right).}
  \label{fig:schematic}
\end{figure*}

The current density-voltage ($J-V$) characteristics under simulated AM 1.5G solar illumination are summarized in Table~\ref{tab:performance}. The baseline cell exhibited a short-circuit current density ($J_{\text{sc}}$) of $14.82\,\text{mA/cm}^2$, an open-circuit voltage ($V_{\text{oc}}$) of $0.58\,\text{V}$, and a fill factor (FF) of $56.2\%$, resulting in a power conversion efficiency (PCE) of $6.81\%$. In contrast, the synapse-engineered cell yielded a significantly higher $J_{\text{sc}}$ of $18.24\,\text{mA/cm}^2$, a $V_{\text{oc}}$ of $0.62\,\text{V}$, and an FF of $61.5\%$. This corresponds to a PCE of $8.42\%$, which constitutes a relative increase of $24\%$ over the baseline device. We also compare our results with a recent high-performance benchmark from Vertex University \cite{sterling2024}.

\begin{table}[h]
  \centering
  \caption{Photovoltaic performance parameters of the baseline and synapse-engineered cells under AM 1.5G illumination ($100\,\text{mW/cm}^2$).}
  \label{tab:performance}
  \small
  \begin{tabularx}{\columnwidth}{@{}l*{4}{C}@{}}
    \toprule
    \textbf{Device Type} & \textbf{$J_{\text{sc}}$} & \textbf{$V_{\text{oc}}$} & \textbf{FF} & \textbf{PCE} \\
    & \textbf{(mA/cm$^2$)} & \textbf{(V)} & \textbf{(\%)} & \textbf{(\%)} \\
    \midrule
    Baseline & 14.82 & 0.58 & 56.2 & 6.81 \\
    Synapse-ETL & 18.24 & 0.62 & 61.5 & 8.42 \\
    Vertex Ref.\ \cite{sterling2024} & 16.10 & 0.59 & 58.0 & 7.15 \\
    \bottomrule
  \end{tabularx}
\end{table}
